% Options for packages loaded elsewhere
\PassOptionsToPackage{unicode}{hyperref}
\PassOptionsToPackage{hyphens}{url}
\PassOptionsToPackage{dvipsnames,svgnames,x11names}{xcolor}
%
\documentclass[
  letterpaper,
  DIV=11,
  numbers=noendperiod]{scrartcl}

\usepackage{amsmath,amssymb}
\usepackage{iftex}
\ifPDFTeX
  \usepackage[T1]{fontenc}
  \usepackage[utf8]{inputenc}
  \usepackage{textcomp} % provide euro and other symbols
\else % if luatex or xetex
  \usepackage{unicode-math}
  \defaultfontfeatures{Scale=MatchLowercase}
  \defaultfontfeatures[\rmfamily]{Ligatures=TeX,Scale=1}
\fi
\usepackage{lmodern}
\ifPDFTeX\else  
    % xetex/luatex font selection
\fi
% Use upquote if available, for straight quotes in verbatim environments
\IfFileExists{upquote.sty}{\usepackage{upquote}}{}
\IfFileExists{microtype.sty}{% use microtype if available
  \usepackage[]{microtype}
  \UseMicrotypeSet[protrusion]{basicmath} % disable protrusion for tt fonts
}{}
\makeatletter
\@ifundefined{KOMAClassName}{% if non-KOMA class
  \IfFileExists{parskip.sty}{%
    \usepackage{parskip}
  }{% else
    \setlength{\parindent}{0pt}
    \setlength{\parskip}{6pt plus 2pt minus 1pt}}
}{% if KOMA class
  \KOMAoptions{parskip=half}}
\makeatother
\usepackage{xcolor}
\setlength{\emergencystretch}{3em} % prevent overfull lines
\setcounter{secnumdepth}{-\maxdimen} % remove section numbering
% Make \paragraph and \subparagraph free-standing
\makeatletter
\ifx\paragraph\undefined\else
  \let\oldparagraph\paragraph
  \renewcommand{\paragraph}{
    \@ifstar
      \xxxParagraphStar
      \xxxParagraphNoStar
  }
  \newcommand{\xxxParagraphStar}[1]{\oldparagraph*{#1}\mbox{}}
  \newcommand{\xxxParagraphNoStar}[1]{\oldparagraph{#1}\mbox{}}
\fi
\ifx\subparagraph\undefined\else
  \let\oldsubparagraph\subparagraph
  \renewcommand{\subparagraph}{
    \@ifstar
      \xxxSubParagraphStar
      \xxxSubParagraphNoStar
  }
  \newcommand{\xxxSubParagraphStar}[1]{\oldsubparagraph*{#1}\mbox{}}
  \newcommand{\xxxSubParagraphNoStar}[1]{\oldsubparagraph{#1}\mbox{}}
\fi
\makeatother


\providecommand{\tightlist}{%
  \setlength{\itemsep}{0pt}\setlength{\parskip}{0pt}}\usepackage{longtable,booktabs,array}
\usepackage{calc} % for calculating minipage widths
% Correct order of tables after \paragraph or \subparagraph
\usepackage{etoolbox}
\makeatletter
\patchcmd\longtable{\par}{\if@noskipsec\mbox{}\fi\par}{}{}
\makeatother
% Allow footnotes in longtable head/foot
\IfFileExists{footnotehyper.sty}{\usepackage{footnotehyper}}{\usepackage{footnote}}
\makesavenoteenv{longtable}
\usepackage{graphicx}
\makeatletter
\def\maxwidth{\ifdim\Gin@nat@width>\linewidth\linewidth\else\Gin@nat@width\fi}
\def\maxheight{\ifdim\Gin@nat@height>\textheight\textheight\else\Gin@nat@height\fi}
\makeatother
% Scale images if necessary, so that they will not overflow the page
% margins by default, and it is still possible to overwrite the defaults
% using explicit options in \includegraphics[width, height, ...]{}
\setkeys{Gin}{width=\maxwidth,height=\maxheight,keepaspectratio}
% Set default figure placement to htbp
\makeatletter
\def\fps@figure{htbp}
\makeatother

\KOMAoption{captions}{tableheading}
\makeatletter
\@ifpackageloaded{caption}{}{\usepackage{caption}}
\AtBeginDocument{%
\ifdefined\contentsname
  \renewcommand*\contentsname{Table of contents}
\else
  \newcommand\contentsname{Table of contents}
\fi
\ifdefined\listfigurename
  \renewcommand*\listfigurename{List of Figures}
\else
  \newcommand\listfigurename{List of Figures}
\fi
\ifdefined\listtablename
  \renewcommand*\listtablename{List of Tables}
\else
  \newcommand\listtablename{List of Tables}
\fi
\ifdefined\figurename
  \renewcommand*\figurename{Figure}
\else
  \newcommand\figurename{Figure}
\fi
\ifdefined\tablename
  \renewcommand*\tablename{Table}
\else
  \newcommand\tablename{Table}
\fi
}
\@ifpackageloaded{float}{}{\usepackage{float}}
\floatstyle{ruled}
\@ifundefined{c@chapter}{\newfloat{codelisting}{h}{lop}}{\newfloat{codelisting}{h}{lop}[chapter]}
\floatname{codelisting}{Listing}
\newcommand*\listoflistings{\listof{codelisting}{List of Listings}}
\makeatother
\makeatletter
\makeatother
\makeatletter
\@ifpackageloaded{caption}{}{\usepackage{caption}}
\@ifpackageloaded{subcaption}{}{\usepackage{subcaption}}
\makeatother

\ifLuaTeX
  \usepackage{selnolig}  % disable illegal ligatures
\fi
\usepackage{bookmark}

\IfFileExists{xurl.sty}{\usepackage{xurl}}{} % add URL line breaks if available
\urlstyle{same} % disable monospaced font for URLs
\hypersetup{
  pdftitle={WCAG 2.1 AA Standards},
  colorlinks=true,
  linkcolor={blue},
  filecolor={Maroon},
  citecolor={Blue},
  urlcolor={Blue},
  pdfcreator={LaTeX via pandoc}}


\title{WCAG 2.1 AA Standards}
\author{}
\date{}

\begin{document}
\maketitle


Includes \textbf{2.1 level AA guidelines only.} These are the guidelines
that need to be met for all content created as part of this project.

\section{Principle 1 - Perceivable}\label{principle-1---perceivable}

Information and user interface components must be presentable to users
in ways they can perceive.

\subsection{Time-Based Media - Guideline
1.2}\label{time-based-media---guideline-1.2}

Provide alternatives for time-based media.

\begin{itemize}
\item
  \textbf{1.2.4 Captions (Live):} Captions are provided for all live
  audio content in synchronized media.
\item
  \textbf{1.2.5 Audio Description (Prerecorded):} Audio description is
  provided for all prerecorded video content in synchronized media.
\end{itemize}

\subsection{Adaptable - Guideline 1.3}\label{adaptable---guideline-1.3}

Create content that can be presented in different ways (for example
simpler layout) without losing information or structure.

\begin{itemize}
\item
  \textbf{1.3.4 Orientation:} Content does not restrict its view and
  operation to a single display orientation, such as portrait or
  landscape, unless a specific display orientation is essential.
\item
  \textbf{1.3.5 Identify Input Purpose:} The purpose of each input field
  collecting information about the user can be programmatically
  determined when:

  \begin{itemize}
  \tightlist
  \item
    The input field serves a purpose identified in the Input Purposes
    for User Interface Components section (link); and
  \item
    The content is implemented using technologies with support for
    identifying the expected meaning for form input data.
  \end{itemize}
\end{itemize}

\subsection{Distinguishable - Guideline
1.4}\label{distinguishable---guideline-1.4}

Make it easier for users to see and hear content including separating
foreground from background.

\begin{itemize}
\item
  \textbf{1.4.3 Contrast (Minimum):} The visual presentation of text and
  images of text has a contrast ratio of at least 4.5:1, except for the
  following:

  \begin{itemize}
  \tightlist
  \item
    \textbf{Large Text:} Large-scale text and images of large-scale text
    have a contrast ratio of at least 3:1;
  \item
    \textbf{Incidental:} Text or images of text that are part of an
    inactive user interface component, that are pure decoration, that
    are not visible to anyone, or that are part of a picture that
    contains significant other visual content, have no contrast
    requirement.
  \item
    \textbf{Logotypes:} Text that is part of a logo or brand name has no
    contrast requirement.
  \end{itemize}
\item
  \textbf{1.4.4 Resize Text:} Except for captions and immages of text,
  text can be resized without assistive technology up to 200 percent
  without loss of content or functionality.
\item
  \textbf{1.4.5 Images of Text:} If the technologies being used can
  achieve the visual presentation, text is used to convey information
  rather than images of text except for the following:

  \begin{itemize}
  \tightlist
  \item
    \textbf{Customizable:} The image of text can be visually customized
    to the user's requirements;
  \item
    \textbf{Essential:} A particular presentation of text is essential
    to the information being conveyed.
  \end{itemize}
\item
  \textbf{1.4.10 Reflow:} Content can be presented without loss of
  information or functionality, and without requiring scrolling in two
  dimensions for:

  \begin{itemize}
  \tightlist
  \item
    Vertical scrolling content at a width equivalent to 320 CSS pixels;
  \item
    Horizontal scrolling content at a height equivalent to 256 CSS
    pixels;
  \end{itemize}

  Except for the parts of the content which require two-dimensional
  layout for usage or meaning.
\item
  \textbf{1.4.11 Non-text Contrast:} The visual presentation of the
  following have a contrast ratio of at least 3:1 against adjacent
  color(s):

  \begin{itemize}
  \tightlist
  \item
    \textbf{User Interface Components:} Visual information required to
    identify user interface components and states, except for inactive
    components or where the appearance of the component is determined by
    the user agent and not modified by the author;
  \item
    \textbf{Graphical Objects:} Parts of graphics required to understand
    the content, except when a particular presentation of graphics is
    essential to the information being conveyed.
  \end{itemize}
\item
  \textbf{1.4.12 Text Spacing:} In content implemented using markup
  languages that support the following text style properties, no loss of
  content or functionality occurs by setting all of the following and by
  changing no other style property:

  \begin{itemize}
  \tightlist
  \item
    Line height (line spacing) to at least 1.5 times the font size;
  \item
    Spacing following paragraphs to at least 2 times the font size;
  \item
    Letter spacing (tracking) to at least 0.12 times the font size;
  \item
    Word spacing to at least 0.16 times the font size.
  \end{itemize}
\item
  \textbf{1.4.13 Content on Hover or Focus:} Where receiving and then
  removing pointer hover or keyboard focus triggers additional content
  to become visible and then hidden, the following are true:

  \begin{itemize}
  \tightlist
  \item
    \textbf{Dismissible:} A mechanism is available to dismiss the
    additional content without moving pointer hover or keyboard focus,
    unless the additional content communicates an input error or does
    not obscure or replace other content;
  \item
    \textbf{Hoverable:} If pointer hover can trigger the additional
    content, then the pointer can be moved over the additional content
    without the additional content disappearing;
  \item
    \textbf{Persistent:} The additional content remains visible until
    the hover or focus trigger is removed, the user dismisses it, or its
    information is no longer valid. Exception: The visual presentation
    of the additional content is controlled by the user agent and is not
    modified by the author.
  \end{itemize}
\end{itemize}

\section{Principle 2 - Operable}\label{principle-2---operable}

User interface components and navigation must be operable.

\subsection{Navigable - Guideline 2.4}\label{navigable---guideline-2.4}

Provide ways to help users navigate, find content, and determine where
they are.

\begin{itemize}
\tightlist
\item
  \textbf{2.4.5 Multiple Ways:} More than one way is available to locate
  a Web page within a set of Web pages except where the Web Page is the
  result of, or a step in, a process.
\item
  \textbf{2.4.6 Headings and Labels:} Headings and labels describe topic
  or purpose.
\item
  \textbf{2.4.7 Focus Visible:} Any keyboard operable user interface has
  a mode of operation where the keyboard focus indicator is visible.
\end{itemize}

\section{Principle 3 -
Understandable}\label{principle-3---understandable}

Information and the operation of the user interface must be
understandable.

\subsection{Readable - Guideline 3.1}\label{readable---guideline-3.1}

Make text content readable and understandable.

\begin{itemize}
\tightlist
\item
  \textbf{3.1.2 Language of Parts:} The human language of each passage
  or phrase in the content can be programmatically determined except for
  proper names, technical terms, words of indeterminate language, and
  words or phrases that have become part of the vernacular of the
  immediately surrounding text.
\end{itemize}

\subsection{Predictable - Guideline
3.2}\label{predictable---guideline-3.2}

Make Web pages appear and operate in predictable ways.

\begin{itemize}
\tightlist
\item
  \textbf{3.2.3 Consistent Navigation:} Navigational mechanisms that are
  repeated on multiple Web pages within a set of Web pages occur in the
  same relative order each time they are repeated, unless a change is
  initiated by the user.
\item
  \textbf{3.2.4 Consistent Identification:} Components that have the
  same functionality within a set of Web pages are identified
  consistently.
\end{itemize}

\subsection{Input Assistance - Guideline
3.3}\label{input-assistance---guideline-3.3}

Help users avoid and correct mistakes.

\begin{itemize}
\item
  \textbf{3.3.3 Error Suggestion:} If an input error is automatically
  detected and suggestions for correction are known, then the
  suggestions are provided to the user, unless it would jeopardize the
  security or purpose of the content.
\item
  \textbf{3.3.4 Error Prevention (Legal, Financial, Data):} For Web
  pages that cause legal commitments or financial transactions for the
  user to occur, that modify or delete user-controllable data in data
  storage systems, or that submit user test responses, at least one of
  the following is true:

  \begin{itemize}
  \tightlist
  \item
    \textbf{Reversible:} Submissions are reversible.
  \item
    \textbf{Checked:} Data entered by the user is checked for input
    errors and the user is provided an opportunity to correct them.
  \item
    \textbf{Confirmed:} A mechanism is available for reviewing,
    confirming, and correcting information before finalizing the
    submission.
  \end{itemize}
\end{itemize}

\section{Principle 4 - Robust}\label{principle-4---robust}

Content must be robust enough that it can be interpreted by a wide
variety of user agents, including assistive technologies.

\subsection{Compatible - Guideline
4.1}\label{compatible---guideline-4.1}

Maximize compatibility with current and future user agents, including
assistive technologies.

\begin{itemize}
\tightlist
\item
  \textbf{4.1.3 Status Message:} In content implemented using markup
  languages, status messages can be programmatically determined through
  role or properties such that they can be presented to the user by
  assistive technologies without receiving focus.
\end{itemize}




\end{document}
